% !TEX root = ../../main.tex

\section{مناقشة وتحليل النتائج} \label{part2:sec4}
% \thispagestyle{empty}

{\fontsize{16}{27}\selectfont
يتناول هذا القسم تحليل البيانات المستخرجة من أداة \el{FFT Analyzer} بعد تطبيق مرشح التمرير المنخفض المصمم، وتقييم جودة الموجات الناتجة مقارنة بالمعايير الهندسية المطلوبة.
\par}

\clearpage

\subsection{تحليل طيف التيار}

بعد تشغيل المحاكاة بوجود المرشح الجديد، أظهرت النتائج تحسناً جذرياً في شكل موجة التيار وتخفيضاً كبيراً في مركباتها التوافقية.

\begin{figure}[H]
    \centering
    \includegraphics[width=0.9\linewidth]{18 - CurrentData with Low-Pass Filter.png}
    \caption{طيف التردد للتيار بوجود مرشح التمرير المنخفض}
    \label{fig:18}
\end{figure}

نلاحظ من الشكل \ref{fig:18} أن بلغت نسبة التشوه التوافقي $THD_I = 4.02\%$. نلاحظ من الرسم البياني اختفاء القمم الحادة عند الترددات التوافقية العالية، حيث عمل المرشح على كبح التيارات التوافقية بنجاح كبير، مما أدى للحصول على موجة تيار تقترب كثيراً من الشكل الجيبي النقي.\clearpage


\subsection{تحليل طيف الجهد}

تعتبر جودة الجهد المعيار الأهم في استقرار الشبكة الكهربائية، وقد أثبت المرشح كفاءة عالية في هذا الجانب.

\begin{figure}[H]
    \centering
    \includegraphics[width=0.9\linewidth]{19 - VoltageData with Low-Pass Filter.png}
    \caption{طيف التردد للجهد بوجود مرشح التمرير المنخفض}
    \label{fig:19}
\end{figure}

حيث انخفضت نسبة التشوه التوافقي لتصل إلى $THD_V = 2.26\%$(كما هو واضح من الشكل \ref{fig:19}). هذه النتيجة تعتبر ممتازة جداً، حيث أنها أقل من $3\%$، وهي الحدود المسموح بها عالمياً حسب النورمات المذكورة في "جلسات العملي".
كما نلاحظ أن منحني الجهد يظهر في السكوب بشكل جيبي منتظم وخالٍ من التشوهات الظاهرة التي كانت موجودة قبل عملية الترشيح.

\subsection{الخلاصة}

يؤكد هذا العمل أن مرشحات التمرير المنخفض \el{(LPF)} هي حلول فعالة واقتصادية لتحسين جودة القدرة وحماية الأجهزة الحساسة من ضجيج التوافقات في الشبكات التي تحتوي على مقومات قدرة. وبالتالي أدى اختيار التردد $150 Hz$ إلى إزاحة نقطة الرنين لأسفل، مما وفر حماية أكبر ضد التوافقية الخامسة وما فوقها. بالإضافة إلى أن الحسابات اليدوية أثبتت دقتها، حيث أدى إدخال القيم ($R, L, C$) المحسوبة إلى استقرار النظام والحصول على نتائج ضمن النطاق المقبول أكاديمياً.
